\section{Synthetic ground truth and the round-trip benchmark}
\label{sec:synthetic-method}

Real scans have no known vertex identities, so correspondence cannot be
measured on them directly. The model can, however, generate its own targets:
sample latent variables, pose the template, export an anonymous mesh, refit
it, and compare against the known truth. The corpus used holds $5000$
specifications drawn with a fixed seed, pose scale $0.50$ and shape scale
$2.0$, with fixed train/test splits; this allows the perturbation, sampling,
initialisation and correspondence variables to be manipulated one at a time.

\begin{algbox}{1: Synthetic round-trip correspondence benchmark}
\begin{enumerate}[nosep,leftmargin=1.4em]
\item Sample $(\bm\beta,\bm\theta)$ from the model's learned distributions (fixed seed).
\item Generate the posed mesh $M$; keep the true identity of every vertex.
\item Optionally corrupt $M$ (noise, holes, dropped points) and strip identities.
\item Run the production fitter on the anonymous mesh, optionally with a
counterfactual input (ground-truth initialisation, regional-identity oracle).
\item For each vertex $i$, compare fitted vertex $i$ with true vertex $i$;
report correspondence, parameter and surface errors per anatomical part.
\end{enumerate}
\end{algbox}

The limitation is intrinsic: the target comes from the model, so the
benchmark bounds error from below on real scans (\S\ref{sec:limitations}).
